% Part: computability
% Chapter: recursive-functions
% Section: partial-functions

\documentclass[../../../include/open-logic-section]{subfiles}

\begin{document}

% SEGMENT OLP-0224-S01

\olfileid{cmp}{rec}{par}
\olsection{பகுதி மீள்வரையறைச் சார்புகள்}


மீள்வரையறைச் சார்புகளின் வரையறைக்கு உள்நோக்கம் தர, முதன்மை
மீள்வரையறைக்குரியவை அல்லாத கணிக்கத்தக்க சார்புகள் உள்ளன என்ற நமது
நிரூபிப்பு உண்மையில் மேலும் வலுவான ஒன்றை நிறுவுகிறது என்பதைக் கவனிக்க.
வாதம் எளியது: $x$ மற்றும் $y$ இன் சார்பாக $f_x(y)$ கணிக்கத்தக்கதாக
இருக்குமாறு சார்புகள் $f_0,f_1,\dots$ ஐ எண்ணிட முடியும் என்ற உண்மையை
மட்டுமே பயன்படுத்தினோம். எனவே இவ்வாதம்
\emph{இவ்வாறு எண்ணிடக்கூடிய எந்தச் சார்பு வகுப்புக்கும்} பொருந்தும்.
இது நம்மை இடர்ப்பாட்டில் வைக்கிறது: கணிக்கத்தக்க சார்புகளை
வெளிப்படையாக விவரிக்க விரும்புகிறோம்; ஆனால் கணிக்கத்தக்க சார்புகளின்
ஒரு திரட்டுக்கான எந்த வெளிப்படையான விளக்கமும் முழுமையானதாக இருக்க
முடியாது!

இதிலிருந்து வெளியேறும் வழி, \emph{பகுதிச்} சார்புகளைப் பயன்படுத்த
அனுமதிப்பதாகும். பகுதியான கணிக்கத்தக்க சார்புகளை எண்ணிடுவது
\emph{சாத்தியம்} என்பதைப் பார்ப்போம்.\iftag{TMs}{ உண்மையில் இது
  அவ்வாறே என்று கிட்டத்தட்ட ஏற்கெனவே அறிந்துள்ளோம்; ஏனெனில் டியூரிங்
  பொறிகளை முறையாக எண்ணிட முடியும்.}{} நமது மூலைவிட்ட வாதத்திற்குப்
பின்னர் திரும்பி, பகுதிச் சார்புகளைச் சேர்க்கும்போது அது ஏன்
செயல்படுவதில்லை என்று ஆராய்வோம்.

இப்போது கேள்வி இதுதான்: அனைத்து பகுதி மீள்வரையறைச் சார்புகளையும்
பெற, முதன்மை மீள்வரையறைச் சார்புகளுடன் எதைச் சேர்க்க வேண்டும்?
இரு செயல்களைச் செய்ய வேண்டும்:
\begin{enumerate}
\item பகுதிச் சார்புகளையும் அனுமதிக்குமாறு முதன்மை மீள்வரையறைச்
  சார்புகளின் வரையறையை மாற்றுக.
\item சில புதிய பகுதிச் சார்புகளும் சேருமாறு வரையறையில் ஏதோ ஒன்றை
  \emph{சேர்க்க}.
\end{enumerate}

முதலாவது எளிது. முன்புபோல் சுழியம், அடுத்தெண், வீழ்த்தல்கள்
ஆகியவற்றில் தொடங்கி, சேர்ப்பு மற்றும் முதன்மை மீள்வரையறையின் கீழ்
மூடுவோம். வரையறையிலுள்ள சில கோவைகள் வரையறுக்கப்படாமல் இருக்கக்கூடிய
சாத்தியத்தை அனுமதிக்க, சேர்ப்பு மற்றும் முதன்மை மீள்வரையறையின்
வரையறைகளை மாற்ற வேண்டும் என்பதே ஒரே வேறுபாடு. $f$ மற்றும் $g$
பகுதிச் சார்புகள் எனில், $f$ ஆனது $x$ இல் வரையறுக்கப்பட்டுள்ளது,
அதாவது $x$ ஆனது $f$ இன் சார்பகத்தில் உள்ளது, என்பதைக் குறிக்க
$f(x) \fdefined$ என எழுதுவோம்; எதிர்மாறாக, $f$ ஆனது~$x$ இல்
வரையறுக்கப்படவில்லை என்பதைக் குறிக்க $f(x) \fundefined$ என எழுதுவோம்.
$f(x)$ மற்றும் $g(x)$ இரண்டும் வரையறுக்கப்படாமல் இருக்கின்றன,
அல்லது இரண்டும் வரையறுக்கப்பட்டுச் சமமாக உள்ளன, என்பதைக் குறிக்க
$f(x) \simeq g(x)$ ஐப் பயன்படுத்துவோம். மேலும் சிக்கலான கோவைகளுக்கும்
இக்குறிமுறைகளைப் பயன்படுத்துவோம். $h$ மற்றும் $g_0$, \dots,~$g_k$
அனைத்தும் பகுதிச் சார்புகள் எனில், பின்வருவதை
\[
h(g_0(\vec x),\dots,g_k(\vec x))
\]
ஒவ்வொரு $g_i$ உம் $\vec x$ இல் வரையறுக்கப்பட்டும், $h$ ஆனது
$g_0(\vec x)$, \dots,~$g_k(\vec x)$ இல் வரையறுக்கப்பட்டும் இருந்தால்,
எனில் மட்டுமே, வரையறுக்கப்பட்டுள்ளது என்ற மரபை ஏற்போம். இப்புரிதலுடன்,
பகுதிச் சார்புகளுக்கான சேர்ப்பு மற்றும் முதன்மை மீள்வரையறையின்
வரையறைகள் மேலே உள்ளவை போன்றவையே; ஆனால் “$=$” என்பதற்குப் பதிலாக
“$\simeq$” ஐ இட வேண்டும்.

பகுதிச் சார்புகளைப் பெறுவதற்கான முதன்மை மீள்வரையறைச் சார்புகளின்
வரையறையில் நாம் சேர்ப்பது \emph{வரம்பற்ற தேடல் செயற்குறி}.
$f(x,\vec z)$ இயல் எண்கள்மீதான ஏதேனும் பகுதிச் சார்பு எனில்,
$\mu x \; f(x,\vec z)$ ஐப் பின்வருமாறு வரையறுக்க:
\begin{quote}
  $f(0,\vec z), f(1,\vec z), \dots, f(x,\vec z)$ அனைத்தும்
  வரையறுக்கப்பட்டும் $f(x,\vec z) = 0$ ஆகவும் இருக்குமாறு
  மிகச்சிறிய $x$ ஒன்று இருந்தால், அந்த $x$
\end{quote}
இல்லையெனில் $\mu x \; f(x,\vec z)$ வரையறுக்கப்படவில்லை என்ற
புரிதலுடன் இதைப் பயன்படுத்துகிறோம். இது $\mu x \; f(x,\vec z)$ ஐத்
தனித்தவாறு வரையறுக்கிறது.

\begin{explain}
நமது வரையறை\iftag{TMs}{ டியூரிங் பொறிகளையோ,}{} வழிமுறைகளையோ,
வேறு எந்தக் குறிப்பிட்ட கணக்கீட்டு மாதிரியையோ சுட்டவில்லை என்பதைக்
கவனிக்க. ஆனால் சேர்ப்பு மற்றும் முதன்மை மீள்வரையறையைப் போலவே,
வரம்பற்ற தேடலுக்குப் பின்னாலும் செயல்முறைசார் கணக்கீட்டு உள்ளுணர்வு
உள்ளது. ஒரு பகுதிச் சார்பின் கணிப்புத்தன்மையைப் பொறுத்தவரை, சார்பு
வரையறுக்கப்படாத உள்ளீடுகள், கணக்கீடு நிற்காத உள்ளீடுகளுக்கு ஒத்தவை.
$\mu x \; f(x,\vec z)$ ஐக் கணிக்கும் செயல்முறை இதுதான்:
0 என்ற மதிப்பு திருப்பித்தரப்படும் வரை $f(0,\vec z), f(1,\vec z),
f(2,\vec z)$ ஆகியவற்றைக் கணிக்க. எனினும் இடைப்பட்ட கணக்கீடுகளில்
ஏதேனும் ஒன்று நிற்காவிட்டால், $\mu x \; f(x,\vec z)$ இன் கணக்கீடும்
நிற்காது.
\end{explain}

$R(x,\vec z)$ ஏதேனும் உறவு எனில், $\mu x \; R(x,\vec z)$ என்பது
$\mu x \; (1 \tsub \Char{R}(x,\vec z))$ என்று வரையறுக்கப்படுகிறது.
வேறுவிதமாகக் கூறினால், $R(x,\vec z)$ பொருந்துமாறு அமையும் மிகச்சிறிய
$x$ மதிப்பை $\mu x \; R(x,\vec z)$ திருப்பித்தருகிறது. எனவே
$f(x,\vec z)$ ஒரு முழுச் சார்பு எனில், $\mu x \; f(x,\vec z)$ ஆனது
$\mu x \; (f(x,\vec z) = 0)$ போன்றதே. ஆனால் நமது மூல வரையறை
மேலும் பொதுவானது என்பதைக் கவனிக்க; ஏனெனில் $f(x,\vec z)$ எல்லா
இடங்களிலும் வரையறுக்கப்படாமல் இருக்கும் சாத்தியத்தை அது அனுமதிக்கிறது
(இதற்கு மாறாக, ஓர் உறவின் சிறப்பியல்புச் சார்பு எப்போதும் முழுச் சார்பு).

\begin{defn}
\emph{பகுதி மீள்வரையறைச் சார்புகளின்} கணம் என்பது, இயல்
எண்களிலிருந்து இயல் எண்களுக்குச் செல்லும் பல்வேறு இடஎண்களுடைய
பகுதிச் சார்புகளின் மிகச்சிறிய கணம்; அது சுழியம், அடுத்தெண்,
வீழ்த்தல்கள் ஆகியவற்றைக் கொண்டு, சேர்ப்பு, முதன்மை மீள்வரையறை,
வரம்பற்ற தேடல் ஆகியவற்றின் கீழ் மூடியது.
\end{defn}

பகுதி மீள்வரையறைச் சார்புகளில் சில, ஒவ்வோர் உள்ளீட்டுக்கும்
வரையறுக்கப்பட்ட முழுச் சார்புகளாகவும் அமையும்.

\begin{defn}
\ollabel{defn:recursive-fn}
\emph{மீள்வரையறைச் சார்புகளின்} கணம் என்பது முழுச் சார்புகளாக
இருக்கும் பகுதி மீள்வரையறைச் சார்புகளின் கணம்.
\end{defn}

ஒரு மீள்வரையறைச் சார்பு எல்லா இடங்களிலும் வரையறுக்கப்பட்டிருப்பதை
வலியுறுத்த, சிலவேளைகளில் “முழு மீள்வரையறைச் சார்பு” எனப்படுகிறது.

\end{document}
